\documentclass[CJK]{beamer}
\usepackage{CJKutf8}
\usepackage{beamerthemesplit}
\usetheme{Malmoe}
\useoutertheme[footline=authortitle]{miniframes}
\usepackage{amsmath}
\usepackage{amssymb}
\usepackage{graphicx}
\usepackage{eufrak}
\usepackage{color}
\usepackage{slashed}
\usepackage{simplewick}
\usepackage{tikz}
\usepackage{tcolorbox}
\usepackage{ulem}
\graphicspath{{../figures/}}
%%figures
\def\lfig#1#2{\includegraphics[width=#1 in]{#2}}
\def\tfig#1#2{\includegraphics[height=#1 in]{#2}}
\def\addfig#1#2{\begin{center}\includegraphics[width=#1 in]{#2}\end{center}}
\def\wulian{\includegraphics[width=0.18in]{emoji_wulian.jpg}}
\def\bigwulian{\includegraphics[width=0.35in]{emoji_wulian.jpg}}
\def\bye{\includegraphics[width=0.18in]{emoji_bye.jpg}}
\def\bigbye{\includegraphics[width=0.35in]{emoji_bye.jpg}}
\def\huaixiao{\includegraphics[width=0.18in]{emoji_huaixiao.jpg}}
\def\bighuaixiao{\includegraphics[width=0.35in]{emoji_huaixiao.jpg}}
\def\jianxiao{\includegraphics[width=0.18in]{emoji_jianxiao.jpg}}
\def\bigjianxiao{\includegraphics[width=0.35in]{emoji_jianxiao.jpg}}
\def\haoqi{\includegraphics[width=0.18in]{emoji_haoqi.jpg}}
%% colors
\def\blacktext#1{{\color{black}#1}}
\def\bluetext#1{{\color{blue}#1}}
\def\redtext#1{{\color{red}#1}}
\def\darkbluetext#1{{\color[rgb]{0,0.2,0.6}#1}}
\def\skybluetext#1{{\color[rgb]{0.2,0.7,1.}#1}}
\def\cyantext#1{{\color[rgb]{0.,0.5,0.5}#1}}
\def\greentext#1{{\color[rgb]{0,0.7,0.1}#1}}
\def\darkgray{\color[rgb]{0.2,0.2,0.2}}
\def\lightgray{\color[rgb]{0.6,0.6,0.6}}
\def\gray{\color[rgb]{0.4,0.4,0.4}}
\def\blue{\color{blue}}
\def\red{\color{red}}
\def\orange{\color[rgb]{1.,0.8,0.}}
\def\green{\color{green}}
\def\darkgreen{\color[rgb]{0,0.4,0.1}}
\def\darkblue{\color[rgb]{0,0.2,0.6}}
\def\skyblue{\color[rgb]{0.2,0.7,1.}}
%%control
\def\uline#1{\underline{\ \hspace{#1 in}\ }}
\def\bral{\left(\begin{array}{l}}
\def\brar{\end{array}\right)}
\def\brall{\left(\begin{array}{ll}}
\def\brarr{\end{array}\right)}
\def\bralll{\left(\begin{array}{lll}}
\def\brarrr{\end{array}\right)}
\def\branchl{\left\{\begin{array}{l}}
\def\branchr{\end{array}\right.}
\def\branchll{\left\{\begin{array}{ll}}
\def\branchrr{\end{array}\right.}
\def\branchlll{\left\{\begin{array}{lll}}
\def\branchrrr{\end{array}\right.}
\def\sfgamma#1{\,\Gamma\left( #1 \right)\,}
\def\be{\begin{equation}}
\def\ee{\nonumber\end{equation}}
\def\bea{\begin{eqnarray}}
\def\eea{\nonumber\end{eqnarray}}
\def\bch{\begin{CJK}{UTF8}{gbsn}}
\def\ech{\end{CJK}}
\def\bchl{\bch\begin{large}}
\def\echl{\end{large}\ech}
\def\bchL{\bch\begin{Large}}
\def\echL{\end{Large}\ech}
\def\bitem{\begin{itemize}}
\def\eitem{\end{itemize}}
\def\bcenter{\begin{center}}
\def\ecenter{\end{center}}
\def\bex{\begin{minipage}{0.2\textwidth}\includegraphics[width=0.6in]{jugelizi.png}\end{minipage}\begin{minipage}{0.76\textwidth}}
\def\eex{\end{minipage}}
\def\chtitle#1{\frametitle{\bch#1\ech}}
\def\bmat#1{\left(\begin{array}{#1}}
\def\emat{\end{array}\right)}
\def\bcase#1{\left\{\begin{array}{#1}}
\def\ecase{\end{array}\right.}
\def\bmini#1{\begin{minipage}{#1\textwidth}}
\def\emini{\end{minipage}}
\def\tbox#1{\begin{tcolorbox}#1\end{tcolorbox}}
\def\pfrac#1#2#3{\left(\frac{\partial #1}{\partial #2}\right)_{#3}}
\def\res#1#2{\,\mathrm{res}\,#1\left(#2\right)\,}
\def\newt#1#2{\left(\begin{array}{c}#1\\ #2\end{array}\right)}
\def\reof#1{\,\mathrm{Re}{\left(#1\right)\,}}
\def\imof#1{\,\mathrm{Im}{\left(#1\right)\,}}
\def\Arg#1{\,\mathrm{Arg}\,#1\,}
%%symbols
%%symbols
\def\bropt{\,(\ \ \ )}
\def\cvmol{C_V^{\rm mol}}
\def\cpmol{C_p^{\rm mol}}
\def\intfull{\int_{-\infty}^\infty}
\def\inthalf{\int_0^\infty}
\def\bropt{\,(\ \ \ )}
\def\optlist#1#2#3{(A)\;#1\ \ \ (B)\;#2\ \ \ (C)\;#3}
\def\ooptlist#1#2#3#4{(A)\;#1\ \ \ (B)\;#2\ \ \ (C)\;#3\ \ (D)\;#4}
\def\sone{$\star$}
\def\stwo{$\star\star$}
\def\sthree{$\star\star\star$}
\def\sfour{$\star\star\star\star$}
\def\sfive{$\star\star\star\star\star$}
\def\rint{{\int_\leftrightarrow}}
\def\roint{{\oint_\leftrightarrow}}
\def\stdHf{{\textit{\r H}_f}}
\def\deltaH{{\Delta \textit{\r H}}}
\def\ii{{\dot{\imath}}}
\def\skipline{{\vskip0.1in}}
\def\skiplines{{\vskip0.2in}}
\def\lagr{{\mathcal{L}}}
\def\hamil{{\mathcal{H}}}
\def\vecv{{\mathbf{v}}}
\def\vecx{{\mathbf{x}}}
\def\vecy{{\mathbf{y}}}
\def\veck{{\mathbf{k}}}
\def\vecp{{\mathbf{p}}}
\def\vecn{{\mathbf{n}}}
\def\vecA{{\mathbf{A}}}
\def\vecP{{\mathbf{P}}}
\def\vecsigma{{\mathbf{\sigma}}}
\def\hatJn{{\hat{J_\vecn}}}
\def\hatJx{{\hat{J_x}}}
\def\hatJy{{\hat{J_y}}}
\def\hatJz{{\hat{J_z}}}
\def\hatj#1{\hat{J_{#1}}}
\def\hatphi{{\hat{\phi}}}
\def\hatq{{\hat{q}}}
\def\hatpi{{\hat{\pi}}}
\def\vel{\upsilon}
\def\Dint{{\mathcal{D}}}
\def\adag{{\hat{a}^\dagger}}
\def\bdag{{\hat{b}^\dagger}}
\def\cdag{{\hat{c}^\dagger}}
\def\ddag{{\hat{d}^\dagger}}
\def\hata{{\hat{a}}}
\def\hatb{{\hat{b}}}
\def\hatc{{\hat{c}}}
\def\hatd{{\hat{d}}}
\def\hatD{{\,\hat{D}}}
\def\hatN{{\hat{N}}}
\def\hatH{{\hat{H}}}
\def\hatp{{\hat{p}}}
\def\Fup{{F^{\mu\nu}}}
\def\Fdown{{F_{\mu\nu}}}
\def\newl{\nonumber \\}
\def\vece{\mathrm{e}}
\def\calM{{\mathcal{M}}}
\def\calT{{\mathcal{T}}}
\def\calR{{\mathcal{R}}}
\def\barpsi{\bar{\psi}}
\def\baru{\bar{u}}
\def\barv{\bar{\upsilon}}
\def\qeq{\stackrel{?}{=}}
\def\ftf{\stackrel{\mathcal{FT}}{\Longrightarrow}}
\def\ftb{\stackrel{\mathcal{FT}}{\Longleftarrow}}
\def\ftfb{\stackrel{\mathcal{FT}}{\Longleftrightarrow}}
\def\ltf{\stackrel{\mathcal{LT}}{\Longrightarrow}}
\def\ltb{\stackrel{\mathcal{LT}}{\Longleftarrow}}
\def\ltfb{\stackrel{\mathcal{LT}}{\Longleftrightarrow}}
\def\torder#1{\mathcal{T}\left(#1\right)}
\def\rorder#1{\mathcal{R}\left(#1\right)}
\def\contr#1#2{\contraction{}{#1}{}{#2}#1#2}
\def\trof#1{\mathrm{Tr}\left(#1\right)}
\def\trace{\mathrm{Tr}}
\def\comm#1{\ \ \ \left(\mathrm{used}\ #1\right)}
\def\tcomm#1{\ \ \ (\text{#1})}
\def\slp{\slashed{p}}
\def\slk{\slashed{k}}
\def\calp{{\mathfrak{p}}}
\def\veccalp{\mathbf{\mathfrak{p}}}
\def\Tthree{T_{\tiny \textcircled{3}}}
\def\pthree{p_{\tiny \textcircled{3}}}
\def\dbar{{\,\mathchar'26\mkern-12mu d}}
\def\erf{\mathrm{erf}}
\def\const{\mathrm{const.}}
\def\pheat{\pfrac p{\ln T}V}
\def\vheat{\pfrac V{\ln T}p}
%%units
\def\fdeg{{^\circ \mathrm{F}}}
\def\cdeg{^\circ \mathrm{C}}
\def\atm{\,\mathrm{atm}}
\def\angstrom{\,\text{\AA}}
\def\SIL{\,\mathrm{L}}
\def\SIkm{\,\mathrm{km}}
\def\SIyr{\,\mathrm{yr}}
\def\SIGyr{\,\mathrm{Gyr}}
\def\SIV{\,\mathrm{V}}
\def\SImV{\,\mathrm{mV}}
\def\SIeV{\,\mathrm{eV}}
\def\SIkeV{\,\mathrm{keV}}
\def\SIMeV{\,\mathrm{MeV}}
\def\SIGeV{\,\mathrm{GeV}}
\def\SIcal{\,\mathrm{cal}}
\def\SIkcal{\,\mathrm{kcal}}
\def\SImol{\,\mathrm{mol}}
\def\SIN{\,\mathrm{N}}
\def\SIHz{\,\mathrm{Hz}}
\def\SIm{\,\mathrm{m}}
\def\SIcm{\,\mathrm{cm}}
\def\SIfm{\,\mathrm{fm}}
\def\SImm{\,\mathrm{mm}}
\def\SInm{\,\mathrm{nm}}
\def\SImum{\,\mathrm{\mu m}}
\def\SIJ{\,\mathrm{J}}
\def\SIW{\,\mathrm{W}}
\def\SIkJ{\,\mathrm{kJ}}
\def\SIs{\,\mathrm{s}}
\def\SIkg{\,\mathrm{kg}}
\def\SIg{\,\mathrm{g}}
\def\SIK{\,\mathrm{K}}
\def\SImmHg{\,\mathrm{mmHg}}
\def\SIPa{\,\mathrm{Pa}}
%page
\def\secpage#1#2{\begin{frame}\bch\bcenter{\bf \Huge #1} \skipline \tbox{#2}\ecenter\ech\end{frame}}

\begin{document}
\bch

\title{第七课  LAPACK，SVD和最佳线性拟合}
\subtitle{http://zhiqihuang.top/algorithms}
  \author{黄志琦}
  \date{}

  \maketitle

  \begin{frame}
    \frametitle{内容摘要}
  \tableofcontents
  \end{frame}

  \section{LAPACK}

  \begin{frame}
    \frametitle{解决线性代数问题的标准库 - LAPACK}
    在正儿八经的大规模科学计算中，线性代数的问题都用{\blue Linear Algebra PACKage (LAPACK)}这个(牛炸天的)程序库来解决。

    \skiplines
    
    大多数同学还在使用windows系统，安装使用LAPACK很麻烦（好吧我承认我跟微软不大对付）。

    \skiplines
    
    所以我们用一个替代品：{\blue python的linalg库}来讲解线性代数的问题。（实际上linalg用的计算内核还是LAPACK。）
  \end{frame}


  \begin{frame}
    \frametitle{Python中的linalg库能干些啥}
    \bitem
  \item{矩阵和矢量的常规操作}
  \item{解线性方程组，求矩阵的逆}
  \item{解决矩阵的本征值问题}
  \item{{\blue 矩阵的SVD}}
  \item{{\blue 解决最佳线性拟合问题}}    
    \eitem

    前面的几种操作线性代数课里都有详细的论述，在linalg里调用也非常容易（请自行搜索练习）。我们要着重讲解的是矩阵的SVD和最佳线性拟合问题。
  \end{frame}
  

  \section{SVD}


  \begin{frame}
    \frametitle{标准化的准对角矩阵}
    如果一个 $m\times n$矩阵$A$的左上最大子方阵（由$A$的$1$至$\min(m, n)$行和$1$至$\min(m, n)$列的公共部分组成）是对角矩阵，且对角元素是逐行递减的非负实数($A_{1,1}\ge A_{2,2}\ge \ldots\ge 0$)，则称$A$是{\blue 标准化的准对角矩阵}。

    \skiplines

    例子：
    
    \bmini{0.25}
    \bcenter
    $3\times 3$，秩为$2$
    \ecenter
   \begin{equation}\left(\begin{array}{lll}
    3 & 0 & 0 \\
    0 & 2 & 0 \\
    0 & 0 & 0 
     \end{array}\right)\nonumber\end{equation}

     \vspace{0.2in}
    \emini    
    \bmini{0.3}
    \bcenter
    $5\times 3$，秩为$3$
    \ecenter
   \begin{equation}\left(\begin{array}{lll}
    3 & 0 & 0 \\
    0 & 1 & 0 \\
    0 & 0 & \frac{1}{2} \\
    0 & 0 & 0 \\
    0 & 0 & 0                         
     \end{array}\right)\nonumber\end{equation}
    \emini
    \bmini{0.38}
    \bcenter
    $4\times 6$，秩为$2$    
    \ecenter
   \begin{equation}\left(\begin{array}{llllll}
     2 & 0 & 0 & 0 & 0 & 0 \\
    0 & 1 & 0 & 0  & 0 & 0 \\
    0 & 0 & 0 & 0  & 0 & 0 \\
    0 & 0 & 0 & 0  & 0 & 0    
     \end{array}\right)\nonumber\end{equation}

     \vspace{0.1in}     
       \emini
    
  \end{frame}


    \begin{frame}
    \frametitle{single value decomposition (SVD)的定义}

    定理：一个$m\times n$的实矩阵$A$，总是能够分解为如下形式：

    $$ A = U \Sigma V^T $$
    其中$U$是$m\times m$的正交矩阵 ($U U^T = I_{m\times m}$)，$V$是$n\times n$的正交矩阵 ($VV^T = I_{n\times n}$)。$\Sigma$是$m\times n$标准化的准对角矩阵。

    （证明略）

    \skipline
    
    这个分解就是名满江湖的single value decomposition ({\blue SVD})。如果要数出5个线性代数最实用理论，SVD一定是其中一个。
    \end{frame}


    \begin{frame}
    \frametitle{SVD的几何含义}
    设$A$是一个$m\times n$的矩阵。映射$\vecx\rightarrow A\vecx$把一个$n$维的“初始空间”里的向量$\vecx$映射到一个$m$维的“目标空间”里的$A\vecx$。

    \skipline
    
    下面我们来通过SVD理解这种不同维空间的映射到底是如何完成的。

    \skipline
    
    设$A$的SVD为：
    $$ A = U \Sigma V^T $$
    $V$的列矢量$\vecv_1, \vecv_2,\ldots, \vecv_n$构成初始空间的一组标准正交基。
 
    $U$的列矢量$\vecu_1, \vecu_2,\ldots, \vecu_n$构成目标空间的一组标准正交基。

    $\Sigma$的准对角元$\sigma_1\ge\sigma_2\ge\ldots\ge\sigma_{\min(m,n)}\ge 0$可以看成一些权重，权重的个数不超过$\min(m,n)$个。非零权重的个数（记为$r$）为$\Sigma$的秩，实际上也是$A$的秩。
    \end{frame}

    \begin{frame}
    \frametitle{SVD的几何含义}    
    如果初始空间的矢量$\vecx$可以按标准正交基分解为
    $$\vecx = VV^T\vecx = V\vecy = y_1 \vecv_1  + y_2 \vecv_2 + \ldots + y_n\vecv_n.$$
    那么
    $$ A \vecx = U\Sigma V^T\vecx = U \Sigma \vecy = (\sigma_1y_1)\vecu_1 + (\sigma_2y_2)\vecu_2+\ldots $$
    也就是说，如果用$V$的列向量为正交基的坐标系描述初始空间，用$U$的列向量为正交基的坐标系描述目标空间，那么映射$A$只是把每个坐标分量依次乘了权重$\sigma_1$, $\sigma_2$，\ldots而已。

    \skipline

    因为在初始空间中取$V$列向量为基的坐标系和在目标空间中取$U$列向量为基的坐标系非常地合理和自然，我们不妨称这些坐标系为“合理坐标系”。


    \end{frame}


    \section{Linear Least Square}
    
    \begin{frame}
      \frametitle{SVD和最佳线性拟合的问题}
      我们经常关心最佳线性拟合问题：对给定的映射$A$以及目标空间里的矢量$\vecb$，如何找到一个初始空间的矢量$\vecx$，使$A\vecx$尽可能地接近$\vecb$。也就是要使$\lVert A\vecx - \vecb\rVert$最小。

      \skipline
      
      有了上述几何观点之后，这个问题现在看起来很容易，不是吗？对$b$的每个“合理坐标系坐标”$b_i$，如果对应的权重$\sigma_i$不为零，则只要取$\vecx$的相应合理坐标系坐标为$b_i/\sigma_i$。对那些权重为零的分量，就可以都取零（当然也可以随便取，不过一般来说我们希望$\vecx$不要包含冗余信息，所以还是取零吧）。如果$\vecb$有维度超出初始空间的分量，就根本不用管了。
      
    \end{frame}


    
    
    



    \begin{frame}
      \frametitle{最佳线性拟合的问题}
      所以求解最佳线性拟合的步骤为：
      \bitem
      \item{SVD: $$A = U\Sigma V^T$$ }
      \item{求$b$的“合理坐标”: $$ \widetilde{\vecb} = U^Tb $$}
      \item{计算$x$的“合理坐标”: $\widetilde{x}_i = \widetilde{b}_i/\sigma_i$ (如果$\sigma_i \ne 0$)或者$\widetilde{x}_i=0$ (如果$\sigma_i=0$).}
      \item{计算$x$的原始坐标系坐标: $\vecx  = V \widetilde{\vecx}$.}
        \eitem

        \skipline
        
        {\bf 注意：}在数值计算时，因为有浮点误差，碰到很小但非零的$\sigma_i$就需要小心：它很可能只是浮点误差产生的，不具有物理意义（具体要根据实际模型来判断），这时需要把它当成零来处理。
    \end{frame}


    \begin{frame}
      \frametitle{linalg的最佳线性拟合函数}
      python的linalg.lstsq集成了上述最佳线性拟合的所有步骤，请自行搜索学习这个函数的用法。在lstsq里当然也允许你手动设置当$\sigma_i$很小时（小于参数rcond)就把它忽略。

      \skipline

      linalg库里也有单独只做SVD的函数linalg.svd。请自行搜索学习用法。
      
    \end{frame}


    \begin{frame}
      \frametitle{补充说明}
      既然python的linalg.lstsq集成了最佳线性拟合的所有步骤，为什么还要唠叨一大堆介绍SVD呢？

      \skiplines
      
      这是因为SVD的几何诠释清晰地告诉你了线性映射$A$做的事情。在实际应用中，很多时候需要先理解模型，再进行数值计算。所以SVD这个中间步骤非常有意义。
      
    \end{frame}

    \section{Applications}
    
    \begin{frame}
      \frametitle{最佳线性拟合的应用}
      家喻户晓的“最小二乘法”的本质是用两个函数
      $$f_1(x) = 1$$
      和
      $$f_2(x) = x$$
      的线性组合来拟合一个和数据符合得最佳的函数 $f(x) $。
    \end{frame}


    \begin{frame}
      \frametitle{最佳线性拟合的应用}
      一般地，可以用任意一组线性独立的函数
      $$f_1(x), \ldots, f_{n}(x).$$
      的线性组合来拟合一个和数据符合得最佳的函数 $f(x)$。


      设数据点有$m$个
      $$(x_1, y_1), (x_2, y_2), \ldots, (x_m, y_m).$$
      设最佳拟合的函数为
      $$ f(x) = c_1 f_1(x) + c_2f_2(x) + \ldots + c_n f_n(x).$$
      用前面学习的套路可以找到$\mathbf{c} = (c_1, c_2,\ldots, c_n)^T$ 令 $\sum_i [y_i - f(x_i)]^2$，也就是$\lVert A\mathbf{c} - \vecy\rVert$最小化。这里$m\times n$的$A$定义为：
      $$A_{i,j} = f_j(x_i)$$
      
    \end{frame}


    \begin{frame}
      \frametitle{作业}
      在$[-\frac{\pi}{2},\frac{\pi}{2}]$上取均匀地取$m=101$个点：
      $$x_0=-\frac{\pi}{2},\ x_1 = -\frac{\pi}{2}+\frac{\pi}{100},\ x_2 = -\frac{\pi}{2}+\frac{2\pi}{100}, \ldots, x_{100} = \frac{\pi}{2}.$$

      制造模拟数据点：
      $$ y_i= \sin(x_i), \ (i=0,1,2,\ldots, 100).$$

      用四次多项式
      $$ y = c_0 + c_1x+c_2x^2+c_3x^3+c_4x^4 $$
      来对这些数据点做最佳线性拟合。把求出的$c_0,c_1,c_2,c_3,c_4$和$\sin x$的泰勒展开系数做比较。
      
    \end{frame}
  \ech
\end{document}
